\subsection{Rang und Lösungstyp}
\label{subsec:linear_systems_solution_criteria}

Mit Hilfe der Zeilenstufenform kann man ablesen, zu welchem Lösungstyp
ein LGS gehört. Eine zentrale Rolle nimmt dabei der Begriff des Rangs ein.

\begin{DefinitionBox}{Rang einer Matrix}{linear_systems_matrix_rank}
  Der \emph{Rang} \EnglishTerm{rank} einer Matrix ist die Anzahl der
  Nichtnullzeilen in einer Zeilenstufenform, wobei eine Nichtnullzeile
  mindestens einen Eintrag ungleich $0$ enthält. Man schreibt
  $\operatorname{rang}(A)$.
  \begin{center}
    \CachedTikzFigure{content/resources/chapters/linear_systems/tikz/row_echelon_rank.tex}
  \end{center}
\end{DefinitionBox}

Auf diese Weise können bei einem LGS sowohl der Rang der Koeffizientenmatrix
$A$ als auch der Rang der erweiterten Koeffizientenmatrix $[A\mid b]$
berechnet werden.

Elementare Zeilenumformungen verändern den Rang nicht. Obwohl eine Matrix auf
verschiedenen Wegen in Zeilenstufenform gebracht werden kann, ist diese Anzahl
daher eindeutig.

\begin{ExampleBox}{Rang ablesen}{linear_systems_rank_example}
  \begin{enumerate}[label=\alph*),leftmargin=1.8em,itemsep=1.3em]
    \item
      \[
        [A\mid b]\sim
        \left[
          \begin{array}{@{}rrr|r@{}}
            \phantom{-}1&\phantom{-}2&\phantom{-}0&\phantom{-}3\\
            \phantom{-}0&\phantom{-}0&\phantom{-}1&\phantom{-}5
          \end{array}
        \right]
        \quad\Longrightarrow\quad
        \begin{aligned}
          n&=\Cloze{1cm}{3},\\[0.5em]
          \operatorname{rang}(A)&=\Cloze{1cm}{2},\\[0.5em]
          \operatorname{rang}([A\mid b])&=\Cloze{1cm}{2}
        \end{aligned}
      \]

    \item
      \[
        [A\mid b]\sim
        \left[
          \begin{array}{@{}rr|r@{}}
            \phantom{-}1&-1&\phantom{-}3\\
            \phantom{-}0&\phantom{-}1&\phantom{-}5\\
            \phantom{-}0&\phantom{-}0&\phantom{-}1
          \end{array}
        \right]
        \quad\Longrightarrow\quad
        \begin{aligned}
          n&=\Cloze{1cm}{2},\\[0.5em]
          \operatorname{rang}(A)&=\Cloze{1cm}{2},\\[0.5em]
          \operatorname{rang}([A\mid b])&=\Cloze{1cm}{3}
        \end{aligned}
      \]
  \end{enumerate}
\end{ExampleBox}

\begin{TheoremBox}{Lösbarkeits- und Rangkriterium}{linear_systems_rank_criterion}
  Für ein LGS mit $n$ Unbekannten, Koeffizientenmatrix $A$ und erweiterter
  Koeffizientenmatrix $[A\mid b]$ gilt:
  \[
    \begin{array}{rcll}
      \operatorname{rang}(A)
      &<&\operatorname{rang}([A\mid b])
      &\Longleftrightarrow\quad \text{keine Lösung},\\[0.7em]
      \operatorname{rang}(A)
      &=&\operatorname{rang}([A\mid b])=n
      &\Longleftrightarrow\quad \text{genau eine Lösung},\\[0.7em]
      \operatorname{rang}(A)
      &=&\operatorname{rang}([A\mid b])<n
      &\Longleftrightarrow\quad \infty\text{-viele Lösungen}.
    \end{array}
  \]
\end{TheoremBox}

Der Rangvergleich prüft zunächst die \emph{Verträglichkeit}
\EnglishTerm{consistency} der Gleichungen. Sind die Ränge gleich, entscheidet
der Vergleich von $\operatorname{rang}(A)$ mit der Zahl $n$ der Unbekannten,
ob genau eine oder $\infty$-viele Lösungen vorliegen.

\clearpage
\begin{ExampleBox}{Lösungstyp an der Zeilenstufenform erkennen}{linear_systems_rank_solution_types_example}
  \begin{enumerate}[label=\alph*),leftmargin=1.8em,itemsep=1.3em]
    \item
      \[
        [A\mid b]\sim
        \left[
          \begin{array}{@{}rr|r@{}}
            1&2&3\\
            0&0&1\\
            0&0&0
          \end{array}
        \right]
        \quad\Longrightarrow\quad
        \left.
        \begin{aligned}
          n&=\Cloze{1cm}{2},\\[0.5em]
          \operatorname{rang}(A)&=\Cloze{1cm}{1},\\[0.5em]
          \operatorname{rang}([A\mid b])&=\Cloze{1cm}{2}
        \end{aligned}
        \quad\right\}
        \quad\Longrightarrow\quad
        \Cloze{4cm}{\text{keine Lösung}}
      \]

    \item
      \[
        [A\mid b]\sim
        \left[
          \begin{array}{@{}rr|r@{}}
            1&2&3\\
            0&1&5
          \end{array}
        \right]
        \quad\Longrightarrow\quad
        \left.
        \begin{aligned}
          n&=\Cloze{1cm}{2},\\[0.5em]
          \operatorname{rang}(A)&=\Cloze{1cm}{2},\\[0.5em]
          \operatorname{rang}([A\mid b])&=\Cloze{1cm}{2}
        \end{aligned}
        \quad\right\}
        \quad\Longrightarrow\quad
        \Cloze{4cm}{\text{genau eine Lösung}}
      \]

    \item
      \[
        [A\mid b]\sim
        \left[
          \begin{array}{@{}rrr|r@{}}
            1&2&0&3\\
            0&0&1&5
          \end{array}
        \right]
        \quad\Longrightarrow\quad
        \left.
        \begin{aligned}
          n&=\Cloze{1cm}{3},\\[0.5em]
          \operatorname{rang}(A)&=\Cloze{1cm}{2},\\[0.5em]
          \operatorname{rang}([A\mid b])&=\Cloze{1cm}{2}
        \end{aligned}
        \quad\right\}
        \quad\Longrightarrow\quad
        \Cloze{4cm}{\infty\text{-viele Lösungen}}
      \]
  \end{enumerate}
\end{ExampleBox}


\begin{KeyPointBox}{Widerspruchszeile}{linear_systems_inconsistent_row}
  Ein lineare Gleich, die niemals erfüllt werden kann, ist
  \[
    0x_1+0x_2+\dots+0x_n=c
    \quad\Longleftrightarrow\quad
    \left[\begin{array}{@{}cccc|c@{}}
      0&0&\cdots&0&c
    \end{array}\right],
    \qquad c\neq0,
  \]
  und heißt \emph{Widerspruchszeile} \EnglishTerm{inconsistent row}. 
  Für ein LGS in Zeilenstufenform gilt daher
  \[
    \exists\;\text{Widerspruchszeile}
    \quad\Longleftrightarrow\quad
    \operatorname{rang}(A)<\operatorname{rang}([A\mid b])
    \quad\Longleftrightarrow\quad
    \text{keine Lösung}.
  \]
\end{KeyPointBox}


\begin{ExampleBox}{Lösungstyp eines $3\times3$-LGS bestimmen}{linear_systems_inconsistent_three_by_three_example}
  Bringen Sie die erweiterte Koeffizientenmatrix des folgenden LGS in
  Zeilenstufenform und bestimmen Sie den Lösungstyp:
  \[
    \begin{aligned}
      x_1+x_2+x_3&=3,\\
      2x_1+x_2-x_3&=1,\\
      3x_1+2x_2&=5
    \end{aligned}
    \quad\Longleftrightarrow\quad
    [A\mid b]=
    \left[
      \begin{array}{@{}rrr|r@{}}
        1&1& 1&3\\
        2&1&-1&1\\
        3&2& 0&5
      \end{array}
    \right].
  \]

  \begin{ClozeArea}[32]{35}
    \textbf{Vorwärtselimination:}
    \[
      \begin{array}{r@{\;}r@{\qquad}l}
        &\left[\begin{array}{@{}rrr|r@{}}
          1&1&1&3\\
          2&1&-1&1\\
          3&2&0&5
        \end{array}\right]
        &\begin{array}{@{}l@{}}
          {}\\
          \smash{\TransformOp{Z_2-2Z_1}}\\
          \smash{\TransformOp{Z_3-3Z_1}}
        \end{array}\\[1.0em]
        \Longleftrightarrow
        &\left[\begin{array}{@{}rrr|r@{}}
          1&1&1&3\\
          0&-1&-3&-5\\
          0&-1&-3&-4
        \end{array}\right]
        &\begin{array}{@{}l@{}}
          {}\\
          {}\\
          \smash{\TransformOp{Z_3-Z_2}}
        \end{array}\\[1.0em]
        \Longleftrightarrow
        &\left[\begin{array}{@{}rrr|r@{}}
          1&1&1&3\\
          0&-1&-3&-5\\
          \textcolor{red!70!black}{0}&\textcolor{red!70!black}{0}&
          \textcolor{red!70!black}{0}&\textcolor{red!70!black}{1}
        \end{array}\right]
        &\text{Zeilenstufenform}
      \end{array}
    \]
    Die letzte Zeile ist die
    \textcolor{red!70!black}{\emph{Widerspruchszeile}}:
    \[
      \textcolor{red!70!black}{
        0x_1+0x_2+0x_3=1
        \quad\Longleftrightarrow\quad
        0=1.
      }
    \]
    \textbf{Rangvergleich:}
    \[
      \left.
      \begin{aligned}
        n&=3,\\[0.5em]
        \operatorname{rang}(A)&=2,\\[0.5em]
        \operatorname{rang}([A\mid b])&=3
      \end{aligned}
      \quad\right\}
      \quad\Longrightarrow\quad
      \text{keine Lösung}.
    \]
    Damit ist die Lösungsmenge
    \[
      \mathcal L=\varnothing.
    \]
  \end{ClozeArea}
\end{ExampleBox}
